\Annex{informative}{Implementation Documentation}{AnnexB}

\Sec{Pre-standard Language Versions}{AnnexB.Versions}

\p HLSL has evolved over time, with various versions of the language being
released alongside different versions of DirectX. Beginning with the DirectX
Shader Compiler (DXC), the compiler began supporting language modes
corresponding to specific historical versions of HLSL. These historical versions
are useful for understanding the evolution of the language and for documenting
observed behaviors in specific implementations, and they are named for the years
in which they were released. The following historical versions of HLSL are
relevant to this specification:

\begin{itemize}
  \item \textbf{HLSL 2015}: Roughly equivalent to the version of HLSL supported
  by the Effects Compiler (FXC).
  \item \textbf{HLSL 2016}: The initial version of HLSL supported by the DirectX
  Shader Compiler (DXC).
  \item \textbf{HLSL 2017}: Added enumerations similar to C++11.
  \item \textbf{HLSL 2018}: Added support for 16-bit data types and intrinsic
  templates.
  \item \textbf{HLSL 2021}: Added user-defined templates, operator overloading,
  short-circuiting boolean operators, and stricter implicit casting rules for
  user-defined structures.
\end{itemize}

\begin{note}
Both the DirectX Shader Compiler (DXC) and Clang support a language mode called
"HLSL 202x" which refers to this draft specification.
\end{note}

\Sec{Known Implementations}{AnnexB.Implementations}

\p This section provides a list of known implementations of HLSL, including both
conforming and non-conforming implementations. This is not meant to be
comprehensive, but is provided to supplement later sections in this annex which
discuss details of specific implementations.

\begin{itemize}
  \item \textbf{Clang}: The open-source C/C++ compiler developed by the LLVM
    project. Clang contains a partial implementation of HLSL this specification
    and is actively being developed. It is available at
    \url{https://clang.llvm.org/}.
  \item \textbf{DirectX Shader Compiler (DXC)}: The open-source reference
    implementation of HLSL, developed by Microsoft and based on the LLVM/Clang
    infrastructure. It is available at
    \url{https://github.com/microsoft/DirectXShaderCompiler}.
  \item \textbf{Effects Compiler (FXC)}: The legacy implementation of HLSL,
    developed by Microsoft. It is included with the Windows SDK and is available
    at \url{https://learn.microsoft.com/en-us/windows/apps/windows-sdk/}.
  \item \textbf{GLSLang}: An open-source implementation of the OpenGL Shading
    Language (GLSL), developed by The Khronos Group. GLSLang contains a (now
    deprecated) partial implementation of HLSL 2018. It is available at
    \url{https://github.com/KhronosGroup/glslang}.
  \item \textbf{Slang}: An open-source shader language and compiler framework
    developed by The Khronos Group. Slang contains a partial implementation of
    HLSL 2018. It is available at \url{https://github.com/shader-slang/slang}.
\end{itemize}

\Sec{Implementation Limits}{AnnexB.Limits}

\p This section describes reasonable implementation limits that may be assumed
by a conforming implementation. These limits specify minimum required or maximum
allowed values for various implementation-defined parameters.

\p \ref{Basic.Types} declares a vector type as having required supported sizes
of 1 to 4 components. Implementations may support larger vector sizes but are
not required to do so. An implementation that supports larger vector sizes must
define its vector size limit.

\Sec{Observed Deviations}{AnnexB.Documentation}

\p This section provides documentation around implementation behaviors that are
known to differ between specific implementations. This includes breaking
changes between historical HLSL and standard HLSL where some compilers or
versions of compilers may not conform to standard HLSL, as well as observed
non-conformant behaviors in such compilers. This is not meant to be a
comprehensive list of bugs.

\Sub{Resource Objects}{AnnexB.Resources}

\p In violation of \ref{Basic.Types.Intangible}, DXC supports storing objects
that contain resources in \texttt{cbuffer} declarations. In this use case the
resource objects are not actually stored in the \texttt{cbuffer}, but instead
the resource is treated as a global resource declaration and name lookup
resolves the value inside the structure declaration to the implied global
resource declaration.

\Sec{Legacy DirectX Features}{AnnexB.LegacyDX}

\p The existing reference compilers support a set of features that worked in
tandem with the Direct3D runtime or other support libraries which are no longer
supported in the latest versions of Direct3D.

\p These features are not included in the standard, nor are they conforming
extensions, as such conforming implementation will not support these features.

\Sub{Effects For Direct3D}{AnnexB.LegacyDX.Effects}

\p The Effects runtime support in Direct3D 9 through Direct3D 11 relied on a set of
syntax extensions for HLSL. The latest documentation for the Direct3D 11 Effects
APIs and shader syntax is available on Microsoft's public documentation
site (\url{https://learn.microsoft.com/en-us/windows/win32/direct3d11/d3d11-graphics-programming-guide-effects}).

\p The main additions to HLSL for the Effects feature are the annotation
syntax (\url{https://learn.microsoft.com/en-us/windows/win32/direct3d11/d3d11-effect-annotation-syntax})
and technique declarations.

\begin{HLSL}
int MyVar <int SomeVal=42;>; // Annotated variable.
sampler S : register(s1) = sampler_state {texture=tex;};

technique10 MyRender // Technique declaration.
{
	pass RenderPass
	{
		SetVertexShader(CompileShader(vs_2_0, VSMain()));
		SetPixelShader(CompileShader(ps_2_0, PSMain()));
	}
}
\end{HLSL}

\p The effects syntax has five keywords: \texttt{sampler\_state},
\texttt{technique}, \texttt{technique10}, \texttt{technique11}, and
\texttt{pass}.

\Sub{Shader Interfaces}{AnnexB.LegacyDX.Interfaces}

\p Direct3D 11 supported a limited level of dynamic symbol resolution based on
\textit{interface declarations}. A shader would declare an interface containing
a set of function declarations, and any class that inherited from that interface
was required to provide definitions of all functions declared in the interface.

\p This functionality used in conjunction with functionality exposed in DXBC and
the Direct3D runtime could swap interface implementations at runtime enabling
code reuse and specialization.

\p This feature is not supported in Direct3D 12, Vulkan, nor any actively
supported shader bytecode. DXC's implementation of the \texttt{interface}
keyword is limited to static verification that objects inheriting from the
interface provide definitions for all interface members.

\p A conforming implementation of this standard will not support the
\textit{interface} keyword.

\Sec{Legacy HLSL Features}{AnnexB.LegacyHLSL}

\p The actively supported reference compiler DXC has support for parsing some
HLSL constructs that had meaning and impact in FXC's implementation, but do
nothing in DXC. These legacy language features are omitted from this standard.

\Sub{cbuffer initializers}{AnnexB.LegacyHLSL.cbuffer}

\p The FXC compiler captured initializers for variables in
\textit{cbuffer-declarations} into shader reflection information. This
information was used by the Effects for Direct3D
(\url{https://github.com/Microsoft/FX11}) runtime, but could also be used
directly by users.

\p The DXC compiler supported parsing initializers but did not capture them into
the shader reflection or any other metadata. This standard disallows
initializing variables in \textit{cbuffer-declarations}, an initializer in such
a declaration is ill-formed and a conforming implementation must issue a
diagnostic.

\Sub{\texttt{uniform} and \texttt{shared} keyword}{AnnexB.LegacyHLSL.Keywords}

\p The current implementation of DXC parses but ignores the \texttt{uniform} and
\texttt{shared} keywords which had meaning in earlier versions of HLSL. A
conforming implementation will not support these keywords, will treat them as
identifier tokens, and will error if they are encountered in grammar formations
where identifiers are disallowed.
